\begin{proof}
En efecto, sea $\mathfrak c$ un ideal de $R$ tal que
$\mathfrak a\cap\mathfrak b\subset\mathfrak c$. Se tiene entonces
$\mathfrak a\cap\mathfrak b\subset\mathfrak a\cap\mathfrak c\subset
\mathfrak c$; ahora bien, $\mathfrak c/(\mathfrak a\cap\mathfrak c)$
es un $R$-módulo isomorfo a
$(\mathfrak a+\mathfrak c)/\mathfrak a$, por tanto a un ideal de
$R/\mathfrak a$, y en consecuencia es de tipo finito. Por otra
parte, $(\mathfrak a\cap\mathfrak c)/(\mathfrak a\cap\mathfrak b)$
es un sub-$R$-módulo de
$\mathfrak a/(\mathfrak a\cap\mathfrak b)$, y este último es
isomorfo a $(\mathfrak a+\mathfrak b)/\mathfrak b$, de modo que es
de tipo finito como ideal de $R/\mathfrak b$. Así,
$(\mathfrak a\cap\mathfrak c)/(\mathfrak a\cap\mathfrak b)$ es
también un $R$-módulo de tipo finito, y finalmente lo mismo vale
para $\mathfrak c/(\mathfrak a\cap\mathfrak b)$.
\end{proof}

